\section*{Trusted quantum samples and reflections: model and promises}
In the $\ApxCount$ problem, our goal is to decide whether a nonempty set $S\subseteq[N]$ satisfies $|S|\ge2w$ (YES) or $|S|\le w$ (NO), promised that one of these is the case. In the query model the algorithm is given a membership oracle for $S$.
QSampling $S$ means that we can approximately prepare the uniform superposition
\[|S\rangle:=\frac1{\sqrt{|S|}}\sum_{i\in S}|i\rangle.\]
To be clear, QSampling supposes that the algorithm is given trusted copies of $|S\rangle$; unlike in the QMA setting, the state need not be ``verified'' by the algorithm.
Reflections about $|S\rangle$ apply the unitary transformation
\[\mathcal R_S:=\id-2|S\rangle\langle S|,\]
which has eigenvalue $-1$ for $|S\rangle$ and eigenvalue $+1$ for all states orthogonal to $|S\rangle$. More concretely, let $U$ be the unitary that performs the map $U|0\rangle=|S\rangle$, for some canonical starting state $|0\rangle$. Since we know the circuit $U$, we can also implement $U^\dagger$, by reversing the order of all the gates and replacing all the gates with their adjoints. Then
\[\mathcal R_S=\id-2|S\rangle\langle S|=U(\id-2|0\rangle\langle0|)U^\dagger.\]
We begin by focusing on the slightly simplified setting where the algorithm is only permitted to perform membership queries, QSamples, and reflections about the state $|S\rangle$. The lower bound below even holds for the exact-size promise $|S|=w$ versus $|S|=2w$; the full unitary model is specified in original Section4.6, retained below.
